\documentclass{article}
\usepackage[margin=1in]{geometry}
\usepackage{amsmath,amsfonts,mathtools,amsthm,amssymb}
\usepackage[l2tabu,orthodox]{nag}
\usepackage{microtype} % fixes spacing or whatever
\usepackage{enumitem}
\usepackage{tikz-cd}
\usepackage{xcolor}
\usepackage{physics}
\usepackage{mathrsfs}
\usepackage{cancel}
% \usepackage{libertine,libertinust1math}
\usepackage{eucal}
\usepackage[toc,page]{appendix}
\usepackage{pgfplots}
\pgfplotsset{compat=1.18}

\newtheorem{proposition}{Proposition}
\newtheorem{theorem}{Theorem}
\newtheorem{lemma}{Lemma}
\newtheorem{corollary}{Corollary}

\theoremstyle{definition}
\newtheorem{remark}{Remark}
\newtheorem{definition}{Definition}
\newtheorem{example}{Example}

% \usepackage[sc,osf]{mathpazo}   % With old-style figures and real smallcaps.
% \linespread{1.025}              % Palatino leads a little more leading
% % Euler for math and numbers
% \usepackage[euler-digits,small]{eulervm}
% \AtBeginDocument{\renewcommand{\hbar}{\hslash}}

\usepackage{etoolbox}
\usepackage{dynkin-diagrams}
\def\row#1/#2!{\dynkin{#1}{#2}\\}
\newcommand{\tble}[1]{
   \renewcommand*\do[1]{\row##1!}
   \[
      \begin{array}{ll}\docsvlist{#1}\end{array}
   \]
}

\allowdisplaybreaks

\renewcommand{\epsilon}{\varepsilon}
% \renewcommand{\phi}{\varphi}

\DeclareMathAlphabet{\mathpzc}{OT1}{pzc}{m}{it}

\newcommand{\goth}[1]{\mathfrak{#1}}
\newcommand{\red}[1]{#1_{\text{red}}}
\newcommand{\ad}[1]{#1_{\text{Ad}}}
\newcommand{\reg}[1]{#1_{\text{reg}}}
\newcommand{\sh}[1]{#1^{\text{sh}}}
\newcommand{\cpdv}[2]{\cfrac{\partial #1}{\partial #2}}

\newcommand{\fA}{\mathcal{A}}
\newcommand{\fB}{\mathcal{B}}
\newcommand{\fC}{\mathcal{C}}
\newcommand{\fD}{\mathcal{D}}
\newcommand{\fE}{\mathcal{E}}
\newcommand{\fF}{\mathcal{F}}
\newcommand{\fG}{\mathcal{G}}
\newcommand{\fH}{\mathcal{H}}
\newcommand{\fI}{\mathcal{I}}
\newcommand{\fJ}{\mathcal{J}}
\newcommand{\fK}{\mathcal{K}}
\newcommand{\fL}{\mathcal{L}}
\newcommand{\fM}{\mathcal{M}}
\newcommand{\fN}{\mathcal{N}}
\newcommand{\fO}{\mathcal{O}}
\newcommand{\fP}{\mathcal{P}}
\newcommand{\fQ}{\mathcal{Q}}
\newcommand{\fR}{\mathcal{R}}
\newcommand{\fS}{\mathcal{S}}
\newcommand{\fT}{\mathcal{T}}
\newcommand{\fX}{\mathcal{X}}
\newcommand{\fY}{\mathcal{Y}}
\newcommand{\fZ}{\mathcal{Z}}

\newcommand{\be}{\mathbf{e}}
\newcommand{\bx}{\mathbf{x}}
\newcommand{\bone}{\mathbf{1}}
\newcommand{\cx}{\bar{x}}
\newcommand{\cy}{\bar{y}}
\newcommand{\vx}{\vec{x}}
\newcommand{\vsigma}{\vec{\sigma}}
\newcommand{\tzeta}{\tilde{\zeta}}

\newcommand{\gm}{\goth{m}}
\newcommand{\gp}{\goth{p}}
\newcommand{\gF}{\goth{F}}
\newcommand{\gU}{\goth{U}}
\newcommand{\gV}{\goth{V}}

\newcommand{\A}{\mathbf{A}}
\newcommand{\C}{\mathbf{C}}
\newcommand{\F}{\mathbf{F}}
\newcommand{\G}{\mathbf{G}}
\newcommand{\bH}{\mathbf{H}}
\newcommand{\N}{\mathbf{N}}
\newcommand{\PP}{\mathbf{P}}
\newcommand{\Q}{\mathbf{Q}}
\newcommand{\R}{\mathbf{R}}
\newcommand{\Z}{\mathbf{Z}}

\newcommand{\dlog}{\dd\log}

\newcommand\srestr[2]{{\left.\kern-\nulldelimiterspace #1\vphantom{\small|} \right|_{#2}}}
\newcommand\restr[2]{{\left.\kern-\nulldelimiterspace #1 \vphantom{\big|} \right|_{#2}}}
\newcommand\gsec[2]{\Gamma({#1}, {#2})}
\newcommand\gsecx[1]{\Gamma(X, {#1})}

\DeclareMathOperator{\chr}{char}
\DeclareMathOperator{\rProj}{\mathbf{Proj}}
\DeclareMathOperator{\rSpec}{\mathbf{Spec}}
\DeclareMathOperator{\rHom}{\mathbf{Hom}}
\DeclareMathOperator{\rExt}{\mathbf{Ext}}
\DeclareMathOperator{\id}{id}
\DeclareMathOperator{\Frac}{Frac}
\DeclareMathOperator{\rk}{rank}
\DeclareMathOperator{\deter}{det}
\DeclareMathOperator{\pic}{Pic}
\DeclareMathOperator{\cacl}{CaCl}
\DeclareMathOperator{\trd}{tr.d.}
\DeclareMathOperator{\cl}{Cl}
\DeclareMathOperator{\depth}{depth}
\DeclareMathOperator{\codim}{codim}
\DeclareMathOperator{\Div}{Div}
\DeclareMathOperator{\coker}{coker}
\DeclareMathOperator{\len}{length}
\DeclareMathOperator{\height}{height}
\DeclareMathOperator{\supp}{Supp}
\DeclareMathOperator{\proj}{Proj}
\DeclareMathOperator{\im}{im}
\DeclareMathOperator{\Hom}{Hom}
\DeclareMathOperator{\Der}{Der}
\DeclareMathOperator{\spec}{Spec}
\DeclareMathOperator{\Aut}{Aut}
\DeclareMathOperator{\ch}{char}
\DeclareMathOperator{\tor}{Tor}
\DeclareMathOperator{\Ann}{Ann}
\DeclareMathOperator{\Syl}{Syl}
\DeclareMathOperator{\Sym}{Sym}
\DeclareMathOperator{\GL}{GL}
\DeclareMathOperator{\SL}{SL}
\DeclareMathOperator{\Stab}{Stab}
\DeclareMathOperator{\Perm}{Perm}
\DeclareMathOperator{\Orb}{Orb}
\DeclareMathOperator{\Gal}{Gal}
\DeclareMathOperator{\Supp}{Supp}
\DeclareMathOperator{\Ext}{Ext}
\DeclareMathOperator{\Tor}{Tor}
\DeclareMathOperator{\PGL}{PGL}
\DeclareMathOperator{\hd}{hd}
\DeclareMathOperator{\pd}{pd}
\DeclareMathOperator{\SO}{SO}
\DeclareMathOperator{\SU}{SU}
\DeclareMathOperator{\Sp}{Sp}
\DeclareMathOperator{\Oth}{O}
\DeclareMathOperator{\Uni}{U}
\DeclareMathOperator{\evf}{ev}
\DeclareMathOperator{\cH}{H}
\DeclareMathOperator{\dR}{R}
\DeclareMathOperator{\End}{End}
\DeclareMathOperator{\Hasse}{Hasse}
\DeclareMathOperator{\Num}{Num}

\author{James Lee}

% Solarized
% \pagecolor[RGB]{0,20,26}
% \color[RGB]{191,191,191}

% Sephia
% \pagecolor[RGB]{249,239,220}

% Grey
% \pagecolor[RGB]{200, 200, 200}

% Black
% \pagecolor[RGB]{0,0,0}
% \color[RGB]{255,255,255}

\title{Chapter 4, Section 3}

\begin{document}
\maketitle
\begin{enumerate} [label=\textbf{\arabic*.}, leftmargin=0em]

\item If $X$ is a curve of genus 2, show that a divisor $D$ is very ample $\iff$ $\deg{D} \geq 5$.
This strengthens (3.3.4).

\begin{proof}
  Let $D$ be a very ample divisor on $X$.
  By the degree-genus formula, a nonsingular plane curve cannot have genus equal to 2, and since $D$ gives an embedding into projective space $|D|$, we must have $\ell(D) = \dim|D| + 1 \geq 4$.
  By Riemann-Roch, we have
  \begin{equation*}
    \deg{D} = \ell(D) - \ell(K - D) + 1,
  \end{equation*}
  so it remains to show $\ell(K - D) = 0$.
  By (Ex. 1.5), $\deg{D} \geq 3$, so $\deg(K - D) < 0$.
  Hence, $\ell(K - D) = 0$.
\end{proof}

\item[\textbf{3.}] If $X$ is a curve of genus $\geq 2$ which is a complete intersection (II, Ex. 8.4) in some $\PP^n$, show that the canonical divisor $K$ is very ample.
Conclude that a curve of genus 2 can never be a complete intersection in any $\PP^n$.

\begin{proof}
  Let $d$ be the degree of $X$ in $\PP^n$.
  Then $X$ is a complete intersection of $n - 1$ hypersurfaces of degree $d_i$, and by (II, Ex. 8.4), we have $\fO_X(K) \cong \fO_X(\ell)$, where $\ell = \sum d_i - n - 1$.
  Thus, we have two ways of computing $\deg{K}$, and $g \geq 2$, so we have
  \begin{equation*}
    \ell d = \deg{K} = 2 g - 2 \geq 2,
  \end{equation*}
  which implies $\ell \geq 2$.
  Hence, $\fO_X(K)$ is very ample.
  If $X$ is a curve of genus 2, then $\ell(K) = 2$, so $K$ corresponds to a morphism $X \to \PP^1$, which can never be an embedding.
\end{proof}

% \item[\textbf{5.}] Let $X$ be a curve in $\PP^3$, which is not contained in any plane.
% \begin{itemize}
%   \item[(a)] If $O \notin X$ is a point, such that the projection from $O$ induces a birational morphism $\varphi$ from $X$ to its image in $\PP^2$, show that $\varphi(X)$ \textit{must} be singular.
%   \item[(b)] If $X$ has degree $d$ and genus $g$, conclude that $g < \frac{1}{2}(d - 1)(d - 2)$.
%   \item[(c)] Now let $\{ X _t\}$ be the flat family of curves induced by the projection (III, 9.8.3) whose fibre over $t = 1$ is $X$, and whose fibre $X_0$ over $t = 0$ is the scheme with support $\varphi(X)$.
%   Show that $X_0$ always has nilpotent elements.
% \end{itemize}

\item[\textbf{6.}] \textit{Curves of Degree 4.}
\begin{itemize}
  \item[(a)] If $X$ is a curve of degree 4 in some $\PP^n$, show that either
  \begin{itemize}
    \item[(1)] $g = 0$, in which case $X$ is either the rational normal quartic in $\PP^4$ (Ex. 3.4) or the rational quartic curve in $\PP^3$ (II, 7.8.6), or
    \item[(2)] $X \subseteq \PP^2$, in which case $g = 3$, or
    \item[(3)] $X \subseteq \PP^3$ and $g = 1$.
  \end{itemize}
  \item[(b)] In the case $g = 1$, show that $X$ is a complete intersection of two irreducible quadric surfaces in $\PP^3$ (I, Ex. 5.1).
\end{itemize}

\begin{proof} $ $ \vspace{0pt}
\begin{itemize} [leftmargin=0cm]
  \item[(a)] Assume $X$ is not contained in any hyperplane of $\PP^n$.
  An embedding $X \to \PP^n$ is given a by linear system $\goth{d}$ of degree $d$ and dimension $n$, where there is a one-to-one correspondence between elements of $\goth{d}$ and hyperplane sections of $X$ in $\PP^n$.
  The linear system $\goth{d}$ corresponds to an invertible sheaf $\fO_X(D)$ on $X$ and a sub-vector space $V \subseteq \Gamma(X, \fO_X(D))$, where $\goth{d} = \PP(V)$.
  By (Ex. 1.5), we have $\dim{\goth{d}} \leq \dim{|D|} \leq \deg{D}$, so we can assume $1 \leq n \leq 4$.

  If $n = 2$, then the degree genus formula gives
  \begin{equation*}
    g = \frac{(d - 1)(d - 2)}{2} = 3.
  \end{equation*}
  If $n = 3$, then $g < 3$ by (III, Ex. 9.11) (Ex. 3.5b).
  We cannot have $g = 2$ by (Ex. 3.1), so $g = 0, 1$, and $X$ is the rational quartic curve in $\PP^3$ if $g = 0$.
  If $n = 4$, $X$ is rational by (Ex. 1.5).

  \item[(b)] Let $\fI$ be the ideal sheaf of $X$ in $\PP^3$, and consider the exact sequence
  \[ \begin{tikzcd}
    0 \arrow[r] & \fI \arrow[r] & \fO_{\PP^3} \arrow[r] & \fO_X \arrow[r] & 0.
  \end{tikzcd} \]
  Twisting by $\fO_{\PP_3}(2)$ and taking cohomology gives
  \[ \begin{tikzcd}
    0 \arrow[r] & \cH^0(\PP^3, \fI(2)) \arrow[r] & \cH^0(\PP^3, \fO_{\PP^3}(2)) \arrow[r] & \cH^0(X, \fO_X(2)) \arrow[r] & \cdots.
  \end{tikzcd} \]
  Now, $X$ is an elliptic curve, and $\fO_X(2)$ corresponds to a divisor $D$ of degree 8.
  Then, Riemann Roch gives $\dim \cH^0(X, \fO_X(2)) = 8$, and the vector space in the middle has dimension 10, so we conclude that $h^0(\PP^3, \fI(2)) \geq 2$.
  An element of $\cH^0(\PP^3, \fI(2))$ is a form of degree 2, whose zero-set will be a surface $Q \subseteq \PP^3$ of degree 2 containing $X$.
  It must be irreducible and reduced because $X$ is not contained in any $\PP^2$, so $X$ must be contained in two distinct irreducible quadric surfaces $Q, Q'$, and $Q \cap Q' = X$ because $X$ has degree 4 (I, 7.7).
\end{itemize} 
\end{proof}
\filbreak


\item[\textbf{9.}] Prove the following lemma of Bertini: if $X$ is a curve of degree $d$ in $\PP^3$, not contained in any plane, then for almost all planes $H \subseteq \PP^3$ (meaning a Zariski open subset of the dual projective space $(\PP^3)^*$), the intersection $X \cap H$ consists of exactly $d$ distinct points, no three of which are collinear.

\begin{proof}
  Let $D$ be a very ample divisor of degree $d$ corresponding to the embedding $X \hookrightarrow \PP^3$.
  Bertini's theorem (II, 8.18) says for all hyperplanes $H$ in $|D|$, the intersection $X \cap H$ consists of exactly $d$ distinct points.
\end{proof}

\item[\textbf{12.}] For each value of $d = 2, 3, 4, 5$ and $r$ satisfying $0 \leq r \leq \frac{1}{2}(d - 1)(d - 2)$, show that there exists an irreducible plane curve of degree $d$ with $r$ nodes and no other singularities.

\begin{proof}
  There is nothing to prove for $d = 2$ or $r = 0$, so assume otherwise.
\end{proof}

\end{enumerate}

\end{document}
